\documentclass[12pt]{scrartcl}
\usepackage[hagavi]{azzam}

% Modulo dan Keterbagian - Materi
% Adaptasi dari 8 Outsource K Math/Pra-OSK/Keterbagian.tex dan Modulo.tex,
% disusun ulang mengikuti pola unit Materi/Latihan/Ujian.
% Gaya penulisan mengikuti mo_material.md: setiap teorema langsung
% diikuti contoh soal, solusi contoh menumpuk persamaan ke bawah.

\title{Modulo dan Keterbagian - Materi}
\date{12 Oktober 2026}

\begin{document}
\maketitle

Modul ini merangkum dua fondasi teori bilangan yang hampir selalu muncul di OSK dan OSP:
keterbagian dan kongruensi modulo. Setiap teorema langsung diikuti contoh soal,
jadi bacalah teori lalu contohnya berpasangan.
Kuasailah sifat dasarnya sampai hafal di luar kepala,
karena hampir setiap soal teori bilangan tingkat lanjut memakainya tanpa menyebutnya lagi.

\section{Keterbagian}

\begin{definition}
Diberikan dua bilangan bulat $a$ dan $b$ dengan $a \ne 0$.
Kita menyebut $a$ \textbf{membagi} $b$ apabila terdapat bilangan bulat $c$
sedemikian hingga $b = ca$. Notasinya adalah $a \mid b$.
\end{definition}

\begin{fact}[Sifat keterbagian]
Untuk bilangan bulat $a, b, c$ dengan pembagi tak nol, berlaku:
\begin{enumerate}
    \item $a \mid 0$ dan $a \mid a$.
    \item Jika $a \mid b$ dan $b \mid c$, maka $a \mid c$.
    \item Jika $a \mid b$ dan $a \mid c$, maka
    $a \mid \alpha b + \beta c$ untuk setiap bilangan bulat $\alpha, \beta$.
    \item Jika $a \mid b$ dengan $b \ne 0$, maka $|a| \le |b|$.
    \item Jika $a \mid b$ dan $b \mid a$, maka $|a| = |b|$.
    \item Jika $a \mid b$, maka $a \mid bc$ untuk setiap bilangan bulat $c$.
\end{enumerate}
\end{fact}

\begin{example}
Carilah semua bilangan bulat $n$ yang memenuhi $n + 1 \mid 100$.
\end{example}
\begin{proof}[Solusi]
Syaratnya berarti $n + 1$ adalah faktor dari $100 = 2^2 \cdot 5^2$.
Faktor positifnya ada $9$ buah, dan masing-masing versi negatifnya juga memenuhi:
\begin{align*}
n + 1 &\in \{\pm 1, \pm 2, \pm 4, \pm 5, \pm 10, \pm 20, \pm 25, \pm 50, \pm 100\}, \\
n &\in \{-101, -51, -26, -21, -11, -6, -5, -3, -2, 0, 1, 3, 4, 9, 19, 24, 49, 99\}.
\end{align*}
\end{proof}

\begin{example}
Carilah semua bilangan bulat $n$ yang memenuhi $n + 1 \mid n^2 + 99$.
\end{example}
\begin{proof}[Solusi]
Kurangkan kelipatan pembagi untuk menyisakan konstanta:
\[ n^2 + 99 = (n + 1)(n - 1) + 100. \]
Karena $n + 1$ selalu membagi $(n + 1)(n - 1)$, syaratnya setara dengan
\[ n + 1 \mid 100. \]
Jawabannya sama seperti contoh sebelumnya.
Trik menyisakan konstanta ini akan muncul berulang kali.
\end{proof}

\begin{theorem}[Algoritma Pembagian]
Untuk setiap bilangan asli $a$ dan $b$, terdapat pasangan \textbf{unik}
bilangan cacah $(q, r)$ yang memenuhi
\[ b = qa + r \quad \text{dengan} \quad 0 \le r < a. \]
\end{theorem}

\begin{example}
Buktikan bahwa di antara $n$ bilangan bulat berurutan, tepat satu habis dibagi $n$.
\end{example}
\begin{proof}[Solusi]
Ambil $n$ bilangan berurutan $m, m + 1, \dots, m + n - 1$ dan bagi $m$ dengan $n$:
\[ m = qn + r \quad \text{dengan} \quad 0 \le r < n. \]
Jika $r = 0$, maka $m$ sendiri yang habis dibagi $n$.
Jika $r > 0$, bilangan ke-$(n - r)$ setelah $m$ adalah
\[ m + (n - r) = qn + r + n - r = (q + 1)n, \]
yang habis dibagi $n$ dan masih di dalam barisan.
Tidak mungkin ada dua, karena selisih dua kelipatan $n$ yang berbeda minimal $n$,
sementara rentang barisan hanya $n - 1$.
\end{proof}

\subsection{Aturan keterbagian bilangan kecil}

\begin{enumerate}
    \item Sebuah bilangan asli habis dibagi $2$ apabila digit terakhirnya genap.
    \item Sebuah bilangan asli habis dibagi $4$ apabila dua digit terakhirnya habis dibagi $4$.
    \item Sebuah bilangan asli habis dibagi $8$ apabila tiga digit terakhirnya habis dibagi $8$.
    \item Secara umum, sebuah bilangan asli habis dibagi $2^n$ apabila
    $n$ digit terakhirnya habis dibagi $2^n$.
    \item Sebuah bilangan asli habis dibagi $5$ apabila digit terakhirnya $0$ atau $5$.
    \item Sebuah bilangan asli habis dibagi $25$ apabila dua digit terakhirnya
    $00$, $25$, $50$, atau $75$.
    \item Secara umum, sebuah bilangan asli habis dibagi $5^n$ apabila
    $n$ digit terakhirnya habis dibagi $5^n$.
    \item Sebuah bilangan asli habis dibagi $3$ maupun $9$ apabila jumlah digitnya
    habis dibagi $3$ maupun $9$.
    \item Sebuah bilangan asli habis dibagi $11$ apabila selisih jumlah digit
    di posisi ganjil dan jumlah digit di posisi genap habis dibagi $11$.
    \item Sebuah bilangan habis dibagi $10^n$ apabila $n$ digit terakhirnya nol.
    \item Sebuah bilangan habis dibagi $ab$ untuk $a$ dan $b$ yang relatif prima
    apabila bilangan tersebut habis dibagi $a$ dan juga habis dibagi $b$.
\end{enumerate}

Untuk bilangan $n$ dengan sedikitnya dua digit, tulis $f(n)$ untuk bilangan yang
diperoleh dengan menghapus digit terakhir $n$, dan $g(n)$ untuk digit terakhir $n$.
Sebagai contoh, $f(421) = 42$ dan $g(421) = 1$. Maka:
\begin{enumerate}
    \item $n$ habis dibagi $7$ apabila $f(n) - 2g(n)$ habis dibagi $7$.
    \item $n$ habis dibagi $13$ apabila $f(n) - 9g(n)$ habis dibagi $13$.
    \item $n$ habis dibagi $17$ apabila $f(n) - 5g(n)$ habis dibagi $17$.
    \item $n$ habis dibagi $19$ apabila $f(n) + 2g(n)$ habis dibagi $19$.
\end{enumerate}

\begin{example}
Buktikan bahwa setiap bilangan $6$ digit berbentuk $\overline{abcabc}$ habis dibagi $11$.
\end{example}
\begin{proof}[Solusi]
Uraikan bilangannya:
\begin{align*}
\overline{abcabc} &= 1000 \cdot \overline{abc} + \overline{abc} \\
&= 1001 \cdot \overline{abc} \\
&= 7 \cdot 11 \cdot 13 \cdot \overline{abc}.
\end{align*}
Jadi bilangan itu bahkan habis dibagi $7$ dan $13$ juga.
\end{proof}

\section{FPB dan KPK}

\begin{definition}
Untuk bilangan asli $m$ dan $n$, $FPB(m, n)$ adalah faktor persekutuan positif
terbesar keduanya, sedangkan $KPK(m, n)$ adalah kelipatan persekutuan positif
terkecil keduanya. Jika $FPB(m, n) = 1$, maka $m$ dan $n$ disebut relatif prima.
\end{definition}

\begin{fact}[Sifat FPB]
\begin{enumerate}
    \item Jika $d = FPB(m, n)$ dan $m = dm'$, $n = dn'$, maka $FPB(m', n') = 1$.
    \item Setiap faktor persekutuan dari $m$ dan $n$ membagi $FPB(m, n)$.
    \item Jika $m = nq + r$, maka $FPB(m, n) = FPB(n, r)$.
    \item Jika $m = p_1^{\alpha_1} \cdots p_k^{\alpha_k}$ dan
    $n = p_1^{\beta_1} \cdots p_k^{\beta_k}$, maka
    \[ FPB(m, n) = p_1^{\min(\alpha_1, \beta_1)} \cdots p_k^{\min(\alpha_k, \beta_k)}. \]
\end{enumerate}
\end{fact}

\begin{theorem}[Algoritma Euclid dan Bezout]
Dengan menerapkan algoritma pembagian berulang,
\[ m = nq_1 + r_1, \quad n = r_1q_2 + r_2, \quad r_1 = r_2q_3 + r_3, \quad \dots \]
sampai sisanya nol, sisa tak nol terakhir adalah $FPB(m, n)$.
Akibatnya, untuk setiap bilangan asli $m$ dan $n$ terdapat bilangan bulat $a$ dan $b$
sedemikian hingga $am + bn = FPB(m, n)$.
\end{theorem}

\begin{example}
Hitunglah $FPB(133, 198)$.
\end{example}
\begin{proof}[Solusi]
Terapkan algoritma Euclid:
\begin{align*}
198 &= 133 + 65, \\
133 &= 2 \cdot 65 + 3, \\
65 &= 21 \cdot 3 + 2, \\
3 &= 2 + 1.
\end{align*}
Sisa tak nol terakhir adalah $1$, jadi
\[ FPB(133, 198) = 1. \]
\end{proof}

\begin{example}
Tentukan bilangan bulat $a$ dan $b$ sehingga $133a + 198b = 1$.
\end{example}
\begin{proof}[Solusi]
Bentuk Bezout diperoleh dengan membaca langkah Euclid di atas mundur dari belakang:
\begin{align*}
1 &= 3 - 2 \\
&= 3 - (65 - 21 \cdot 3) \\
&= 22 \cdot 3 - 65 \\
&= 22 \cdot (133 - 2 \cdot 65) - 65 \\
&= 22 \cdot 133 - 45 \cdot 65 \\
&= 22 \cdot 133 - 45 \cdot (198 - 133) \\
&= 67 \cdot 133 - 45 \cdot 198.
\end{align*}
Jadi $a = 67$ dan $b = -45$ memenuhi.
\end{proof}

\begin{fact}[Sifat KPK]
\begin{enumerate}
    \item Setiap kelipatan persekutuan dari $m$ dan $n$ adalah kelipatan dari $KPK(m, n)$.
    \item Dengan notasi faktorisasi prima yang sama,
    \[ KPK(m, n) = p_1^{\max(\alpha_1, \beta_1)} \cdots p_k^{\max(\alpha_k, \beta_k)}. \]
    \item $FPB(m, n) \times KPK(m, n) = mn$.
\end{enumerate}
\end{fact}

\begin{example}
Hitunglah $KPK(24, 36)$ dan periksa $FPB(24, 36) \times KPK(24, 36) = 24 \cdot 36$.
\end{example}
\begin{proof}[Solusi]
Faktorkan kedua bilangan:
\begin{align*}
24 &= 2^3 \cdot 3, \\
36 &= 2^2 \cdot 3^2.
\end{align*}
Ambil pangkat terkecil untuk FPB dan terbesar untuk KPK:
\begin{align*}
FPB(24, 36) &= 2^2 \cdot 3 = 12, \\
KPK(24, 36) &= 2^3 \cdot 3^2 = 72.
\end{align*}
Pemeriksaan:
\[ 12 \cdot 72 = 864 = 24 \cdot 36. \]
\end{proof}

\section{Kongruensi Modulo}

\begin{definition}
Misalkan $a$, $b$, dan $n$ bilangan bulat dengan $n \ne 0$.
Kita menyebut $a$ \textbf{kongruen} dengan $b$ modulo $n$ apabila $n \mid a - b$.
Notasinya adalah $a \equiv b \pmod{n}$.
\end{definition}

\begin{fact}[Sifat kongruensi]
\begin{enumerate}
    \item $a \equiv a \pmod{n}$.
    \item Jika $a \equiv b \pmod{n}$, maka $b \equiv a \pmod{n}$.
    \item Jika $a \equiv b \pmod{n}$ dan $b \equiv c \pmod{n}$, maka $a \equiv c \pmod{n}$.
    \item Jika $a \equiv b \pmod{n}$, maka $ka \equiv kb \pmod{n}$ untuk setiap bilangan bulat $k$.
    \item Jika $a \equiv b \pmod{n}$ dan $c \equiv d \pmod{n}$, maka
    $ja + kc \equiv jb + kd \pmod{n}$ untuk setiap bilangan bulat $j$ dan $k$.
    \item Jika $a \equiv b \pmod{n}$ dan $c \equiv d \pmod{n}$, maka $ac \equiv bd \pmod{n}$.
    \item $a \equiv b \pmod{n_i}$ untuk $i = 1, \dots, k$ jika dan hanya jika
    $a \equiv b \pmod{KPK(n_1, \dots, n_k)}$.
\end{enumerate}
\end{fact}

\begin{example}
Tentukan sisa pembagian $2^{10} + 3^{10}$ oleh $7$.
\end{example}
\begin{proof}[Solusi]
Sederhanakan masing-masing suku modulo $7$:
\begin{align*}
2^3 = 8 &\equiv 1 \pmod{7}, \\
2^{10} = (2^3)^3 \cdot 2 &\equiv 1^3 \cdot 2 \equiv 2 \pmod{7}.
\end{align*}
\begin{align*}
3^2 = 9 &\equiv 2 \pmod{7}, \\
3^4 &\equiv 2^2 \equiv 4 \pmod{7}, \\
3^8 &\equiv 4^2 = 16 \equiv 2 \pmod{7}, \\
3^{10} = 3^8 \cdot 3^2 &\equiv 2 \cdot 2 \equiv 4 \pmod{7}.
\end{align*}
Jumlahkan:
\[ 2^{10} + 3^{10} \equiv 2 + 4 \equiv 6 \pmod{7}. \]
Sisanya $6$.
\end{proof}

\subsection{Teorema penting}

\begin{theorem}[Invers modular]
Untuk setiap bilangan asli $m$ dan $n$ yang relatif prima,
terdapat bilangan asli $x$ sedemikian hingga $mx \equiv 1 \pmod{n}$.
\end{theorem}

\begin{example}
Tentukan bilangan asli $x$ sehingga $7x \equiv 1 \pmod{31}$.
\end{example}
\begin{proof}[Solusi]
Terapkan Euclid pada $31$ dan $7$:
\begin{align*}
31 &= 4 \cdot 7 + 3, \\
7 &= 2 \cdot 3 + 1.
\end{align*}
Baca mundur untuk bentuk Bezout:
\begin{align*}
1 &= 7 - 2 \cdot 3 \\
&= 7 - 2 \cdot (31 - 4 \cdot 7) \\
&= 9 \cdot 7 - 2 \cdot 31.
\end{align*}
Modulo $31$ suku terakhir lenyap:
\[ 9 \cdot 7 \equiv 1 \pmod{31}. \]
Jadi $x = 9$ memenuhi (periksa: $63 = 2 \cdot 31 + 1$).
\end{proof}

\begin{theorem}[Euler]
Misalkan $\varphi(m)$ menyatakan banyak bilangan asli tidak lebih dari $m$
yang relatif prima dengan $m$. Jika faktorisasi prima dari $n$ adalah
\[ n = p_1^{\alpha_1} p_2^{\alpha_2} \cdots p_k^{\alpha_k}, \]
maka
\[ \varphi(n) = n\left(1 - \frac{1}{p_1}\right)\left(1 - \frac{1}{p_2}\right)
\cdots\left(1 - \frac{1}{p_k}\right). \]
Untuk $a$ dan $m$ yang relatif prima, berlaku $a^{\varphi(m)} \equiv 1 \pmod{m}$.
\end{theorem}

\begin{example}
Hitunglah $\varphi(100)$, lalu tentukan sisa $3^{41}$ dibagi $100$.
\end{example}
\begin{proof}[Solusi]
Faktorisasi $100 = 2^2 \cdot 5^2$ memberi
\[ \varphi(100) = 100\left(1 - \frac{1}{2}\right)\left(1 - \frac{1}{5}\right) = 100 \cdot \frac{1}{2} \cdot \frac{4}{5} = 40. \]
Karena $3$ relatif prima dengan $100$, teorema Euler memberi
\[ 3^{40} \equiv 1 \pmod{100}. \]
Maka
\[ 3^{41} = 3 \cdot 3^{40} \equiv 3 \cdot 1 \equiv 3 \pmod{100}. \]
Sisanya $3$.
\end{proof}

\begin{theorem}[Fermat Kecil]
Untuk setiap bilangan prima $p$ dan bilangan bulat $a$, berlaku
$a^p \equiv a \pmod{p}$. Jika $p \nmid a$, maka $a^{p-1} \equiv 1 \pmod{p}$.
\end{theorem}

\begin{example}
Carilah sisa pembagian $2^{100}$ oleh $13$.
\end{example}
\begin{proof}[Solusi]
Karena $13$ prima dan tidak membagi $2$, Fermat Kecil memberi
\[ 2^{12} \equiv 1 \pmod{13}. \]
Bagi pangkatnya:
\[ 100 = 12 \cdot 8 + 4. \]
Maka
\begin{align*}
2^{100} &= (2^{12})^8 \cdot 2^4 \\
&\equiv 1^8 \cdot 16 \\
&\equiv 3 \pmod{13}.
\end{align*}
Sisanya $3$.
\end{proof}

\begin{theorem}[Wilson]
Untuk setiap bilangan prima $p$, berlaku $(p-1)! \equiv -1 \pmod{p}$.
\end{theorem}

\begin{example}
Tentukan sisa pembagian $14!$ oleh $17$.
\end{example}
\begin{proof}[Solusi]
Wilson memberi
\[ 16! \equiv -1 \pmod{17}. \]
Uraikan $16!$:
\begin{align*}
16! &= 16 \cdot 15 \cdot 14! \\
&\equiv (-1) \cdot (-2) \cdot 14! \\
&\equiv 2 \cdot 14! \pmod{17}.
\end{align*}
Jadi
\[ 2 \cdot 14! \equiv -1 \pmod{17}. \]
Karena $2^{-1} \equiv 9 \pmod{17}$ (sebab $2 \cdot 9 = 18 \equiv 1$):
\[ 14! \equiv -1 \cdot 9 \equiv -9 \equiv 8 \pmod{17}. \]
Sisanya $8$.
\end{proof}

\begin{theorem}[Sisa Tiongkok / CRT]
Misalkan $m_1, \dots, m_k$ bilangan asli lebih dari $1$ yang saling relatif prima.
Maka untuk setiap bilangan bulat $a_1, \dots, a_k$, sistem kekongruenan
$x \equiv a_i \pmod{m_i}$ memiliki solusi yang tunggal modulo $m_1 m_2 \cdots m_k$.
\end{theorem}

\begin{example}
Tentukan bilangan asli terkecil $x$ dengan $x \equiv 2 \pmod{3}$ dan $x \equiv 3 \pmod{5}$.
\end{example}
\begin{proof}[Solusi]
Dari syarat pertama, $x = 3k + 2$ untuk suatu bilangan cacah $k$.
Substitusikan ke syarat kedua:
\begin{align*}
3k + 2 &\equiv 3 \pmod{5}, \\
3k &\equiv 1 \pmod{5}.
\end{align*}
Karena $3 \cdot 2 = 6 \equiv 1 \pmod{5}$, nilai $k$ terkecil adalah $k = 2$:
\[ x = 3 \cdot 2 + 2 = 8. \]
Periksa: $8 \equiv 2 \pmod{3}$ dan $8 \equiv 3 \pmod{5}$.
\end{proof}

\begin{example}[AIME 1983]
Carilah sisa pembagian $6^{83} + 8^{83}$ oleh $49$.
\end{example}
\begin{proof}[Solusi]
Tulis $6 = 7 - 1$ dan uraikan dengan binomial.
Semua suku dengan $7^2$ atau lebih tinggi habis dibagi $49$, jadi
\begin{align*}
6^{83} = (7-1)^{83} &\equiv (-1)^{83} + 83 \cdot 7 \cdot (-1)^{82} \\
&\equiv -1 + 581 \\
&\equiv 41 \pmod{49}.
\end{align*}
Dengan cara yang sama:
\begin{align*}
8^{83} = (7+1)^{83} &\equiv 1 + 83 \cdot 7 \\
&\equiv 43 \pmod{49}.
\end{align*}
Jumlahkan:
\[ 41 + 43 = 84 \equiv 35 \pmod{49}. \]
\end{proof}

\section{Bilangan Prima}

\begin{definition}
Bilangan asli $p > 1$ disebut \textbf{prima} apabila satu-satunya pembagi
positifnya adalah $1$ dan $p$ sendiri. Bilangan asli lebih dari $1$ yang bukan
prima disebut \textbf{komposit}.
\end{definition}

\begin{theorem}[Faktorisasi prima]
Setiap bilangan asli $n > 1$ dapat dinyatakan secara unik sebagai perkalian
bilangan prima:
\[ n = p_1^{\alpha_1} p_2^{\alpha_2} \cdots p_k^{\alpha_k} \]
dengan $p_1, \dots, p_k$ prima berbeda dan $\alpha_i$ bilangan asli.
\end{theorem}

\begin{example}
Tentukan banyak faktor positif dari $2004$.
\end{example}
\begin{proof}[Solusi]
Faktorkan dahulu:
\[ 2004 = 4 \cdot 501 = 2^2 \cdot 3 \cdot 167. \]
Setiap faktor positif berbentuk $2^a 3^b 167^c$ dengan
$0 \le a \le 2$, $0 \le b \le 1$, $0 \le c \le 1$.
Banyak pilihannya:
\[ (2 + 1)(1 + 1)(1 + 1) = 3 \cdot 2 \cdot 2 = 12. \]
\end{proof}

\begin{fact}[Sifat yang sering dipakai]
\begin{enumerate}
    \item Banyaknya bilangan prima tak hingga.
    \item Setiap bilangan prima lebih dari $3$ berbentuk $6k \pm 1$ untuk suatu bilangan asli $k$.
    \item Jika $p$ prima dan $p \mid ab$, maka $p \mid a$ atau $p \mid b$ (atau keduanya).
\end{enumerate}
\end{fact}

\begin{example}
Buktikan bahwa $24 \mid p^2 - 1$ untuk setiap bilangan prima $p > 3$.
\end{example}
\begin{proof}[Solusi]
Karena $p$ ganjil, $p - 1$ dan $p + 1$ dua bilangan genap berurutan,
jadi salah satunya habis dibagi $4$:
\[ 8 \mid (p-1)(p+1) = p^2 - 1. \]
Di sisi lain $p \equiv \pm 1 \pmod{3}$, sehingga
\begin{align*}
p^2 &\equiv 1 \pmod{3}, \\
3 &\mid p^2 - 1.
\end{align*}
Karena $8$ dan $3$ relatif prima:
\[ 24 \mid p^2 - 1. \]
\end{proof}

\section{Residu Pangkat}

Tidak semua sisa bisa muncul sebagai nilai suatu pangkat.
Keterbatasan ini sering menjadi kunci untuk menunjukkan suatu persamaan tidak punya solusi.

\begin{fact}[Residu kuadrat]
\begin{enumerate}
    \item Sisa $x^2$ dibagi $3$ hanya bisa $0$ atau $1$.
    \item Sisa $x^2$ dibagi $4$ hanya bisa $0$ atau $1$.
    \item Sisa $x^2$ dibagi $5$ hanya bisa $0$, $1$, atau $4$.
    \item Sisa $x^2$ dibagi $8$ hanya bisa $0$, $1$, atau $4$.
\end{enumerate}
\end{fact}

\begin{example}
Buktikan tidak ada bilangan asli $a$, $b$, $c$ yang memenuhi $a^2 + b^2 - 8c = 6$.
\end{example}
\begin{proof}[Solusi]
Tinjau persamaan itu modulo $8$:
\[ a^2 + b^2 \equiv 6 \pmod{8}. \]
Sisa kuadrat modulo $8$ hanya $0$, $1$, atau $4$,
sehingga jumlah dua sisa hanya bisa
\[ 0,\ 1,\ 2,\ 4,\ \text{atau}\ 5 \pmod{8}. \]
Nilai $6$ tidak pernah muncul, jadi persamaan itu tidak punya solusi.
\end{proof}

\begin{fact}[Residu pangkat lebih tinggi]
\begin{enumerate}
    \item Sisa $x^3$ dibagi $7$ hanya bisa $0$, $1$, atau $6$.
    \item Sisa $x^3$ dibagi $9$ hanya bisa $0$, $1$, atau $8$.
    \item Sisa $x^4$ dibagi $16$ hanya bisa $0$ atau $1$.
\end{enumerate}
\end{fact}

\begin{example}
Buktikan tidak ada bilangan bulat $x$ sehingga $x^3 \equiv 2 \pmod{7}$.
\end{example}
\begin{proof}[Solusi]
Hitung semua kemungkinan $x^3$ modulo $7$:
\begin{align*}
0^3 &\equiv 0 \pmod{7}, \\
1^3 &\equiv 1 \pmod{7}, \\
2^3 = 8 &\equiv 1 \pmod{7}, \\
3^3 = 27 &\equiv 6 \pmod{7}, \\
4^3 = 64 &\equiv 1 \pmod{7}, \\
5^3 = 125 &\equiv 6 \pmod{7}, \\
6^3 = 216 &\equiv 6 \pmod{7}.
\end{align*}
Sisa yang muncul hanya $0$, $1$, dan $6$.
Nilai $2$ tidak pernah muncul.
\end{proof}

\begin{example}
Carilah semua bilangan cacah $n$ sehingga $n^2 + 100$ merupakan bilangan kuadrat.
\end{example}
\begin{proof}[Solusi]
Misalkan $n^2 + 100 = m^2$ untuk suatu bilangan cacah $m$.
Pindahkan ruas dan faktorkan:
\[ (m - n)(m + n) = 100, \qquad m + n \ge m - n \ge 0. \]
Pasangan faktor $(m-n, m+n)$ yang jumlahnya genap hanyalah
\[ (2, 50) \quad \text{dan} \quad (10, 10), \]
yang memberi
\[ (m, n) = (26, 24) \quad \text{dan} \quad (10, 0). \]
Jadi $n = 0$ atau $n = 24$.
\end{proof}

\section{Referensi}
\begin{itemize}
    \item Modul Pra-OSK: Keterbagian (Azzam dan Akeyla).
    \item Modul Pra-OSK: Modulo (Azzam dan Akeyla).
\end{itemize}

\end{document}
